\documentclass[11pt]{article}
\usepackage[utf8]{inputenc}
\usepackage[T1]{fontenc}
\usepackage{amsmath}
\usepackage{amsfonts}
\usepackage{amssymb}
\usepackage[version=4]{mhchem}
\usepackage{stmaryrd}

\begin{document}

\section{Optimal Subsidies in the Quality Ladders Model}

Take the quality ladders model as studied in lectures, with some modifications. The physical marginal cost of making a machine of variety $\nu$ at quality $q$ is given by $\psi q(\nu,t \mid q)$. However, the government pays a subsidy $S_M \in (0,1)$ such that the actual marginal cost faced by machine producers is $\psi_e q(\nu,t \mid q)$, where $\psi_e \equiv (1-S_M)\psi.$

The government also pays a subsidy $S_{R\&D} \in (0,1)$ for R\&D expenditure. (Assume that the government budget is balanced by imposing lump-sum taxes on consumers.) Let us also set $\theta = 1$.\\

1. The final-good production function is $Y(t)=\frac{1}{1-\beta}L^\beta \int_0^1 q(\nu,t)x(\nu,t)^{1-\beta}\, d\nu .$ Show that the demand curve for a given variety is $x(\nu,t \mid q)=\left(\frac{q(\nu,t)}{p(\nu,t \mid q)}\right)^{1/\beta}L.$\\

2. The owner of the patent for variety $\nu$ at quality $q$ maximizes profits $\pi(\nu,t \mid q)=p(\nu,t \mid q)x(\nu,t \mid q)-\psi_e q(\nu,t \mid q)x(\nu,t \mid q).$ Imposing the normalization $1-\beta \equiv \psi$ from now on, show that $p(\nu,t \mid q)=(1-S_M)q(\nu,t)$ and $\pi(\nu,t \mid q)=(1-S_M)^{\frac{\beta-1}{\beta}}\beta q(\nu,t)L.$\\


3.Let $V(\nu,t \mid q)$ be the value of a patent that a successful innovator obtains if she raises the quality of variety $\nu$ to the level $q(\nu,t)$. Spending one unit of the final good on research generates a flow rate of innovation $\eta/q^I$, where $q^I=\frac{q(\nu,t)}{\lambda}$ is the incumbent's quality level. Interpret the condition for free entry into research, $1-S_{R\&D}=\frac{\eta V(\nu,t \mid q)}{q(\nu,t)/\lambda}.$\\

4. Interpret the no-arbitrage condition $r(t)V(\nu,t \mid q)=\pi(\nu,t \mid q)-z(\nu,t \mid q)V(\nu,t \mid q)+\dot V(\nu,t \mid q).$\\

5. Explain why under free entry, $\dot V(\nu,t \mid q)=0.$ Use the free entry and the no-arbitrage conditions, and the expression for profits you derived to show that $r(t)+z(\nu,t \mid q)=s\lambda\eta\beta L,$ where $s\equiv\frac{(1-S_M)^{\frac{\beta-1}{\beta}}}{1-S_{R\&D}}.$\\

6. On a balanced growth path (BGP), the following equations must hold: $r+z=s\lambda\eta\beta L,$ $g=r-\rho,$ and $g=(\lambda-1)z.$ Show that $g=\frac{s\lambda\eta\beta L-\rho}{(\lambda-1)^{-1}+1}.$\\

7. The socially optimal growth rate is $g^*=(\lambda-1)(1-\beta)^{-1/\beta}\eta\beta L-\rho.$ Show that to achieve this growth rate, it must be that $s=(1-\beta)^{-1/\beta}-\frac{\rho}{(\lambda-1)\lambda\eta\beta L}.$\\

8. True or false? Explain.
\begin{quote}
``Any combination of $S_M$ and $S_{R\&D}$ that satisfies equation $s=(1-\beta)^{-1/\beta}-\frac{\rho}{(\lambda-1)\lambda\eta\beta L}$ will achieve the socially optimal allocation.''
\end{quote}






\end{document}